\chapter{State, Dynamics, and Estimation}

\section*{학습 목표}
Chapter 1은 Real VLA를 다섯 개의 책임 layer로 정의했다. Chapter 2는 그 중 observation과 state의 차이를 기술적으로 고정한다. 이 장의 목표는 다음과 같다.

\begin{enumerate}
  \item observation $\obs_t$와 physical state $\state_t$의 차이를 설명한다.
  \item VLA policy, estimator, controller가 서로 다른 시간축과 정보 집합을 사용한다는 점을 명확히 한다.
  \item partial observability, latency, calibration error가 왜 VLA failure mode가 되는지 정식화한다.
  \item Chapter 3의 control interface와 Chapter 13의 action representation 논의를 위한 state notation을 준비한다.
\end{enumerate}

\section{로봇이 본다는 것은 상태를 안다는 뜻이 아니다}

VLA 논문에서 흔히 ``image와 instruction을 입력으로 action을 출력한다''고 말한다. 이 표현은 모델 구조를 간단히 설명할 때는 유용하지만, 실제 로봇 시스템을 설명하기에는 부족하다. 카메라 image는 observation이다. 로봇이 제어해야 하는 것은 physical state이다. 두 개는 같지 않다.

예를 들어 머그컵이 보인다고 해서 로봇이 머그컵의 6D pose, 표면 마찰, 내부 액체 유무, gripper와의 접촉 상태, 유리컵과의 최소 거리, 카메라 timestamp와 controller timestamp의 오차를 모두 알고 있다는 뜻은 아니다. VLM이 장면을 잘 설명해도 그 설명은 상태 추정을 대체하지 않는다.

\begin{definitionbox}{Observation과 state}
Observation $\obs_t$는 센서가 시점 $t$에서 제공한 측정값이다. State $\state_t$는 로봇과 환경의 물리 상태이다. 대부분의 실제 로봇 문제에서 $\state_t$는 직접 관측되지 않으며, 시스템은 estimate $\hat{\state}_t$ 또는 belief $\bel_t$를 사용한다.
\end{definitionbox}

Real VLA에서 이 구분이 중요한 이유는 단순하다. VLA policy는 주로 observation-conditioned policy로 학습되지만, controller와 safety layer는 state-dependent constraint를 요구한다. 즉 policy가 보는 것과 controller가 필요로 하는 것이 다르다.

\section{State-space formulation}

가장 기본적인 로봇 시스템 모델은 다음 state-space form으로 쓸 수 있다 \cite{kalman1960,thrun2005}.

\begin{align}
  \state_{t+1} &= f(\state_t, \ctrlcmd_t, \bm{w}_t), \label{eq:state-dynamics}\\
  \obs_t &= h(\state_t, \bm{v}_t). \label{eq:observation-model}
\end{align}

여기서 $f$는 dynamics, $h$는 observation model, $\bm{w}_t$는 process noise, $\bm{v}_t$는 measurement noise이다. VLA policy가 직접 사용하는 입력은 보통 $\obs_t$ 또는 observation history $\obs_{\leq t}$이다. 그러나 실제 dynamics는 $\state_t$와 control command $\ctrlcmd_t$에 의해 전개된다.

Manipulator의 경우 state는 다음처럼 구성될 수 있다.
\begin{equation}
  \state_t =
  \left[
  \conf_t^\top,
  \dot{\conf}_t^\top,
  \bm{x}_{ee,t}^\top,
  \bm{x}_{obj,t}^\top,
  \bm{c}_t^\top
  \right]^\top ,
  \label{eq:manipulator-state}
\end{equation}
여기서 $\conf_t$는 joint configuration, $\dot{\conf}_t$는 joint velocity, $\bm{x}_{ee,t}$는 end-effector pose, $\bm{x}_{obj,t}$는 object pose, $\bm{c}_t$는 contact-related variable이다. 모든 항이 항상 직접 관측되는 것은 아니다.

VLA policy는 이 전체 state 대신 다음과 같은 입력을 받을 수 있다.
\begin{equation}
  \act_t = \policy(\bm{I}_t, \bm{p}_t, \task, m_t),
  \label{eq:image-policy}
\end{equation}
여기서 $\bm{I}_t$는 image observation, $\bm{p}_t$는 proprioception, $\task$는 language instruction, $m_t$는 memory이다. 이 식은 \cref{eq:state-dynamics}와 \cref{eq:observation-model}을 제거하지 않는다. 단지 policy가 사용하는 정보 집합을 나타낼 뿐이다.

\section{Partial observability}

로봇 조작은 대부분 partially observable problem이다. 관측되지 않는 변수는 많다.

\begin{itemize}
  \item 물체의 정확한 6D pose와 covariance
  \item 접촉 여부와 접촉 force
  \item gripper slip
  \item 가려진 obstacle
  \item 표면 마찰과 object compliance
  \item actuator delay와 backlash
  \item 사람의 의도 또는 workspace 침입 가능성
\end{itemize}

따라서 Real VLA 시스템은 observation을 곧바로 state로 취급하면 안 된다. 더 일반적으로는 belief state를 사용한다.

\begin{equation}
  \bel_t = p(\state_t \mid \obs_{1:t}, \ctrlcmd_{1:t-1}, \task).
  \label{eq:belief-definition}
\end{equation}

Belief update는 다음과 같은 recursive form으로 쓸 수 있다.

\begin{align}
  \bar{\bel}_t(\state_t) &=
  \int p(\state_t \mid \state_{t-1}, \ctrlcmd_{t-1})
  \bel_{t-1}(\state_{t-1})\,d\state_{t-1}, \label{eq:belief-predict}\\
  \bel_t(\state_t) &=
  \eta\,p(\obs_t \mid \state_t)\bar{\bel}_t(\state_t). \label{eq:belief-correct}
\end{align}

여기서 \cref{eq:belief-predict}는 prediction step이고, \cref{eq:belief-correct}는 correction step이다. EKF, UKF, particle filter, visual-inertial odometry, learned state estimator는 이 구조의 다른 구현으로 볼 수 있다.

\section{VLA policy가 belief를 사용하는 방식}

VLA policy가 명시적으로 $\bel_t$를 입력으로 받지 않더라도, 시스템 수준에서는 belief가 존재한다. 예를 들어 다음 세 가지 구현이 가능하다.

\begin{enumerate}
  \item \textbf{Implicit belief}: policy가 image history와 proprioception history를 직접 받아 내부 hidden state로 불확실성을 흡수한다.
  \item \textbf{Explicit estimator}: 별도 state estimator가 $\hat{\state}_t$ 또는 covariance $\Sigma_t$를 계산하고, policy 또는 controller가 이를 사용한다.
  \item \textbf{Hybrid interface}: VLA policy는 image와 language를 사용하고, controller/safety layer는 estimator가 만든 geometric state와 uncertainty를 사용한다.
\end{enumerate}

Real VLA 책에서 중요한 것은 어떤 방식이 항상 옳다는 결론이 아니다. 중요한 것은 interface를 명시하는 것이다. Policy가 uncertainty를 직접 다루는가? Controller가 별도 state estimate를 받는가? Safety layer가 covariance를 볼 수 있는가? Evaluation log에 estimator failure가 기록되는가?

\begin{table}[ht]
\centering
\caption{State information이 stack에서 사용되는 위치}
\label{tab:state-stack}
\small
\begin{tabularx}{\textwidth}{p{30mm}X X}
\toprule
Layer & 필요한 state 정보 & state error가 만드는 failure \\
\midrule
Embodied reasoning & object identity, relation, approximate pose & 잘못된 object reference 또는 infeasible subgoal \\
VLA policy & image, proprioception, short history, estimated target & action drift, wrong grasp, layout shift failure \\
Motion/control & joint state, end-effector pose, velocity, constraint & tracking error, collision, unstable motion \\
Safety/assurance & forbidden region, distance, force, velocity, uncertainty & late stop, unsafe recovery, false accept \\
Data/evaluation & logged state, timestamp, intervention, failure label & benchmark artifact, hidden estimator failure \\
\bottomrule
\end{tabularx}
\end{table}

\section{Dynamics and controller-relevant state}

Controller가 필요로 하는 state는 policy가 보는 observation보다 더 엄격하다. Manipulator dynamics는 일반적으로 다음과 같이 쓸 수 있다.

\begin{equation}
  \bm{M}(\conf)\ddot{\conf}
  + \bm{C}(\conf,\dot{\conf})\dot{\conf}
  + \bm{g}(\conf)
  + \bm{\tau}_{contact}
  = \bm{\tau}.
  \label{eq:manipulator-dynamics}
\end{equation}

여기서 $\bm{M}$은 inertia matrix, $\bm{C}$는 Coriolis/centrifugal term, $\bm{g}$는 gravity, $\bm{\tau}_{contact}$는 contact force가 joint torque로 반영된 항, $\bm{\tau}$는 actuator torque이다.

VLA policy가 waypoint나 delta pose를 출력하더라도, 실제 robot은 \cref{eq:manipulator-dynamics}의 제약 아래 움직인다. 따라서 controller는 적어도 다음 정보를 필요로 한다.

\begin{itemize}
  \item 현재 joint configuration $\conf_t$와 velocity $\dot{\conf}_t$
  \item end-effector pose와 Jacobian
  \item obstacle 또는 workspace constraint
  \item velocity, acceleration, torque limit
  \item contact state 또는 force estimate
\end{itemize}

이 정보가 잘못되면 VLA policy가 semantic하게 맞는 action을 내도 물리 실행은 실패할 수 있다.

\section{Latency and time alignment}

VLA 시스템에서 time index는 자주 무시되지만, 실제 deployment에서는 핵심 변수이다. Camera frame, proprioception, VLA inference, safety monitor, controller loop는 서로 다른 rate로 동작한다.

\begin{table}[ht]
\centering
\caption{Real VLA stack의 대표 time scale}
\label{tab:timescale}
\small
\begin{tabularx}{\textwidth}{p{34mm}X X}
\toprule
Component & 대표 rate 또는 지연 & 영향 \\
\midrule
Camera / vision encoder & 10--60 Hz, frame latency 존재 & object pose와 실제 pose 사이 time skew 발생 \\
VLA inference & model 크기에 따라 수 Hz 이하일 수 있음 & high-rate correction이 어려움 \\
Controller loop & 100 Hz--1 kHz & fast stabilization과 contact response 담당 \\
Safety monitor & event에 따라 high-rate 필요 & 늦으면 stop이 의미 없어짐 \\
Data logger & multi-rate synchronization 필요 & failure analysis와 reproducibility에 영향 \\
\bottomrule
\end{tabularx}
\end{table}

측정 시각과 실행 시각이 다르면 policy는 오래된 observation에 대해 action을 낼 수 있다. 이를 단순하게 쓰면
\begin{equation}
  \act_t = \policy(\obs_{t-\delta_o}, \conf_{t-\delta_q}, \task),
  \label{eq:delayed-policy}
\end{equation}
이다. 여기서 $\delta_o$와 $\delta_q$는 observation과 proprioception의 effective delay이다. Controller가 실제로 받는 명령은 다시 inference delay와 communication delay를 거친다.
\begin{equation}
  \ctrlcmd_t = \controller(\act_{t-\delta_\pi}, \hat{\state}_{t-\delta_s}, \mathcal{C}_{t-\delta_c}).
  \label{eq:delayed-control}
\end{equation}

이 식은 추상적이지만 중요한 engineering message를 준다. VLA가 ``현재 장면''을 보고 있다고 말하려면 timestamp alignment를 증명해야 한다. 특히 contact-rich manipulation과 human-proximate robot에서는 delay가 safety failure로 직접 이어질 수 있다.

\section{Calibration and coordinate frames}

State estimation은 센서 값을 추정 state로 바꾸는 문제일 뿐 아니라 coordinate frame을 맞추는 문제이기도 하다. VLA action이 end-effector delta인지, camera frame waypoint인지, world frame pose인지에 따라 controller interpretation이 달라진다.

대표적으로 camera frame의 point $\bm{x}^{C}$를 robot base frame으로 바꾸려면
\begin{equation}
  \bm{x}^{B} = \bm{T}_{BC}\bm{x}^{C}
  \label{eq:frame-transform}
\end{equation}
와 같은 extrinsic transform이 필요하다. 여기서 $\bm{T}_{BC}\in SE(3)$는 camera-to-base transform이다. Calibration error가 있으면 policy가 맞는 object를 선택해도 controller는 잘못된 위치로 움직인다.

Frame mismatch는 VLA action representation과 직접 연결된다.

\begin{itemize}
  \item Action token은 frame 정보를 token vocabulary에 암묵적으로 숨길 수 있다.
  \item Continuous delta는 scale과 frame convention을 명시해야 한다.
  \item Waypoint는 world frame, base frame, camera frame 중 어느 frame인지 명확해야 한다.
  \item Whole-body action은 base motion과 arm motion의 frame coupling을 다룬다.
\end{itemize}

따라서 Chapter 13의 action representation 비교는 Chapter 2의 state/frame notation 위에 세워진다.

\section{State uncertainty and safety}

Safety layer는 단순히 estimated state만 보면 부족하다. Uncertainty도 봐야 한다. 물체 위치 추정치가 같아도 covariance가 큰 경우와 작은 경우는 다른 safety decision을 요구한다.

장애물과 end-effector 사이 signed distance를 $d(\hat{\state}_t)$라 쓰자. Uncertainty를 반영한 conservative safety margin은 다음처럼 둘 수 있다.
\begin{equation}
  d(\hat{\state}_t) - \lambda \sqrt{\bm{J}_d \Sigma_t \bm{J}_d^\top} \ge d_{\min},
  \label{eq:uncertainty-safety}
\end{equation}
여기서 $\Sigma_t$는 state covariance, $\bm{J}_d$는 distance function의 local Jacobian, $\lambda$는 confidence margin, $d_{\min}$은 최소 안전 거리이다.

이 식은 safety가 단순 rule이 아니라 state estimation과 연결된다는 점을 보여준다. VLA policy가 semantic하게 ``glass를 피하라''를 이해해도, glass pose uncertainty가 크면 안전한 motion을 계산할 수 없다.

\section{Estimation failure modes}

Real VLA에서 estimation failure는 model failure처럼 보일 수 있다. 그러나 원인을 분리해야 한다.

\begin{table}[ht]
\centering
\caption{Estimation-related failure mode}
\label{tab:estimation-failure}
\small
\begin{tabularx}{\textwidth}{p{32mm}X X}
\toprule
Failure mode & 증상 & 확인할 log \\
\midrule
Object pose error & grasp가 물체 옆을 지나감 & pose estimate, image timestamp, calibration \\
State latency & moving object를 늦게 따라감 & sensor/control timestamp, inference latency \\
Frame mismatch & 일정 방향으로 systematic offset 발생 & camera-base extrinsic, action frame convention \\
Contact unobservability & grasp 후 slip을 감지하지 못함 & tactile/force sensor, gripper current, contact label \\
Hidden intervention & success처럼 보이나 사람이 중간에 수정 & intervention log, reset policy, video audit \\
\bottomrule
\end{tabularx}
\end{table}

이 표는 Chapter 16의 evaluation protocol과 연결된다. Benchmark가 final success만 기록하면 estimation failure가 policy success 또는 policy failure로 잘못 분류될 수 있다.

\begin{casebox}{Case study: VLA failure처럼 보이는 calibration failure}
카메라 extrinsic이 몇 cm 어긋나면 VLA는 target object를 맞게 찾고도 gripper가 항상 물체 옆을 지나갈 수 있다. 이 실패는 language grounding 실패도, action head 실패도 아니다. Real VLA report는 camera-base calibration, timestamp, frame convention, pose covariance를 함께 기록해야 한다.
\end{casebox}

\chapterwork
{Kalman filtering과 Probabilistic Robotics의 belief-state 관점을 읽고, VLA observation pipeline이 어떤 state estimate를 암묵적으로 요구하는지 확인하라 \cite{kalman1960,thrun2005}.}
{image-conditioned policy가 있다고 해서 state estimation 문제가 사라지지 않는 이유는 무엇인가?}
{mug/rack/glass task에서 $\obs_t$, $\hat{\state}_t$, $\bel_t$, calibration metadata, safety constraint를 분리한 state record를 설계하라.}

\section{Roboticist's Takeaway}

\begin{definitionbox}{Chapter 2 Takeaway}
VLA가 observation을 입력으로 받는다는 사실은 state estimation 문제가 사라졌다는 뜻이 아니다. Real VLA system은 $\obs_t$, $\hat{\state}_t$, $\bel_t$, $\act_t$, $\ctrlcmd_t$가 어떤 시간축과 frame에서 연결되는지 명시해야 한다. 이 interface가 불명확하면 policy, controller, safety, evaluation failure를 구분할 수 없다.
\end{definitionbox}

다음 장은 이 상태 추정과 시간 정렬 위에서 VLA output이 어떤 controller interface를 통과해야 하는지 다룬다.
